\documentclass[11pt]{article}
\usepackage[a4paper,margin=28mm]{geometry}
\usepackage{amsmath,amssymb}
\usepackage{hyperref}
\setlength{\emergencystretch}{2em}
\title{A nonintegral metrized graph with rational canonical Green values}
\author{Ziang Ni (\texttt{ziangni-sys})}
\date{10 October 2026}
\begin{document}
\maketitle
\noindent\textbf{Disproof of conjecture 00000008370.}
We refute the universal transcendence assertion for nonintegral
real-weight graphs by constructing the continuous Arakelov--Green
function on a metrized segment.

Let $\Gamma=[0,r]$ with its ordinary path metric, for $r>0$. It is a
connected metrized graph consisting of one edge of length $r$. Its
effective resistance is $\rho(x,y)=|x-y|$. With the standard graph
Laplacian convention
\[
 \Delta f=-f''\,dx-\sum_p\Big(\sum_{v\text{ outgoing at }p}D_vf(p)\Big)\delta_p,
\]
the derivative jump of $|x-y|$ at an interior $y$ is $2$, while the
endpoint terms are $\delta_0+\delta_r$. Thus
\[
 \Delta_x\rho(x,y)=\delta_0+\delta_r-2\delta_y.
\]
The same formula holds at $y=0,r$ after combining coincident atoms.
The canonical measure, defined by
$\mu_{\mathrm{can}}=\tfrac12\Delta_x\rho(x,y)+\delta_y$, is therefore
\[
 \mu_{\mathrm{can}}=\tfrac12(\delta_0+\delta_r).
\]
This is the continuous metrized-graph canonical measure, not a
replacement by a finite-matrix normalization. For the standard
Laplacian and canonical measure definitions, see Baker--Faber,
\emph{Metrized Graphs, Laplacian Operators, and Electrical Networks},
\S3 and \S7, especially Definition 9 and the interval example preceding
Theorem 11; available at
\href{https://arxiv.org/abs/math/0407428}{arXiv:math/0407428}.

Define the full kernel, for all $x,y\in\Gamma$, by
\[
 g(x,y)=\frac r4-\frac12|x-y|.
\]
It is continuous and symmetric. Fix an interior point $y$.
The $x$-derivatives are $1/2$ for $x<y$ and $-1/2$ for $x>y$.
The negative derivative jump is $1$, and the negative outward endpoint
slopes give $-\delta_0/2-\delta_r/2$. Consequently
\[
 \Delta_xg(x,y)=\delta_y-\mu_{\mathrm{can}}.
\]
For $y=0$ the constant edge derivative is $-1/2$, giving
$\Delta g=(\delta_0-\delta_r)/2$; for $y=r$ it is $1/2$, giving
the opposite measure. These are exactly the same Green equation.
Furthermore, for every $0\leq y\leq r$,
\[
 \int_\Gamma g(x,y)\,d\mu_{\mathrm{can}}(x)
 =\frac12\left(\frac r4-\frac y2+\frac r4-\frac{r-y}{2}\right)=0.
\]
Thus this kernel has precisely the defining Green equation and
normalization of the canonical Arakelov--Green function. Uniqueness
also follows directly: the difference of two such functions has zero
Laplacian, hence zero edge second derivative and zero endpoint slopes;
it is constant, and zero canonical mean makes that constant zero.

Take $r=1/2$. The edge length is positive and nonintegral, while
\[
 g(0,0)=\frac18\ne0.
\]
This is rational and therefore algebraic, contrary to universal
transcendence on nonintegral real-weight graphs. The other conjuncts
are irrelevant once this assertion fails. This example does not use
zero Green values, a vertex-only surrogate, or an assumed Green
certificate.

\paragraph{Lean coverage.}
The project defines the full real kernel on the interval and proves
its actual derivatives on both affine pieces. It computes endpoint
fluxes from those derivatives and verifies the distributional graph
Laplacian against arbitrary test functions, including both endpoint
cases, as well as its canonical mean zero. It proves positivity and
nonintegrality of the length and the nonzero algebraic diagonal value.
The finite atomic signed measure is represented by its action on test
functions; the continuous second-derivative contribution is zero
because the proven derivatives are constant on each piece.
Lean 4.19.0 and Mathlib are pinned, with printed standard-axiom audits.
\end{document}
